% ============================================================================
\clearpage
\section{Secuencia de construcción}\label{sec:secuencia}
% ============================================================================

La construcción avanza por etapas con un criterio verificable de salida (\tabref{tab:secuencia}). Ninguna
etapa se considera terminada porque su código exista: se considera terminada cuando supera su criterio.

\begin{table}[H]
\centering\small
\caption{Etapas, entregables y criterios para avanzar.}
\label{tab:secuencia}
\begin{tabularx}{\linewidth}{@{}>{\RaggedRight}p{2.6cm} Y Y@{}}
\toprule
\encab{Etapa} & \encab{Entregable} & \encab{Criterio para avanzar}\\
\midrule
Datos y especificación & Muestra auditada, eventos, horarios y contrato de datos & Cadenas reconstruibles en
el instante de decisión; cobertura y costo conocidos\\
Distribuciones & Superficies restringidas, cuantiles, bandas y \emph{replay} & Coherencia matemática y
estabilidad frente a spreads y colas\\
Dinámica & Transporte, factores filtrados e innovaciones & Separar ruido de cambios persistentes; cálculo
sin información futura\\
Pronóstico & Referencias y modelos $\P$ con calibración temporal & Mejora incremental fuera de muestra o
delimitación honesta de su utilidad\\
Decisión & Escenarios, CVaR, costos y zona de no operación & Riesgo factible y resultados estables en
pruebas de estrés\\
Paper & Informe diario y reconciliación & Operación reproducible y fallos manejados sin órdenes duplicadas\\
\bottomrule
\end{tabularx}
\end{table}

La estimación inicial de ingeniería es de 6 a 10 semanas para un prototipo integrado con un universo
pequeño, si los datos adecuados están disponibles (\figref{fig:calendario_obra}). No incluye esperar
suficientes observaciones prospectivas para evaluar habilidad predictiva, y se revisa después de auditar la
muestra de datos.

\grafica[0.62]{s9_calendario}{Secuencia tentativa. El reparto por etapa es ilustrativo y se superpone
parcialmente; lo que la propuesta fija es el total de 6 a 10 semanas de ingeniería, que se revisará tras la
auditoría de datos.}{fig:calendario_obra}

\begin{noconcluir}
La investigación puede concluir que el transporte ayuda al riesgo y no al rendimiento, o que no añade valor
sobre las referencias. Ambos son resultados válidos del programa, no fracasos del calendario.
\end{noconcluir}
